% Options for packages loaded elsewhere
\PassOptionsToPackage{unicode}{hyperref}
\PassOptionsToPackage{hyphens}{url}
%
\documentclass[
]{article}
\usepackage{amsmath,amssymb}
\usepackage{iftex}
\ifPDFTeX
  \usepackage[T1]{fontenc}
  \usepackage[utf8]{inputenc}
  \usepackage{textcomp} % provide euro and other symbols
\else % if luatex or xetex
  \usepackage{unicode-math} % this also loads fontspec
  \defaultfontfeatures{Scale=MatchLowercase}
  \defaultfontfeatures[\rmfamily]{Ligatures=TeX,Scale=1}
\fi
\usepackage{lmodern}
\ifPDFTeX\else
  % xetex/luatex font selection
\fi
% Use upquote if available, for straight quotes in verbatim environments
\IfFileExists{upquote.sty}{\usepackage{upquote}}{}
\IfFileExists{microtype.sty}{% use microtype if available
  \usepackage[]{microtype}
  \UseMicrotypeSet[protrusion]{basicmath} % disable protrusion for tt fonts
}{}
\makeatletter
\@ifundefined{KOMAClassName}{% if non-KOMA class
  \IfFileExists{parskip.sty}{%
    \usepackage{parskip}
  }{% else
    \setlength{\parindent}{0pt}
    \setlength{\parskip}{6pt plus 2pt minus 1pt}}
}{% if KOMA class
  \KOMAoptions{parskip=half}}
\makeatother
\usepackage{xcolor}
\usepackage{longtable,booktabs,array}
\usepackage{calc} % for calculating minipage widths
% Correct order of tables after \paragraph or \subparagraph
\usepackage{etoolbox}
\makeatletter
\patchcmd\longtable{\par}{\if@noskipsec\mbox{}\fi\par}{}{}
\makeatother
% Allow footnotes in longtable head/foot
\IfFileExists{footnotehyper.sty}{\usepackage{footnotehyper}}{\usepackage{footnote}}
\makesavenoteenv{longtable}
\usepackage{graphicx}
\makeatletter
\def\maxwidth{\ifdim\Gin@nat@width>\linewidth\linewidth\else\Gin@nat@width\fi}
\def\maxheight{\ifdim\Gin@nat@height>\textheight\textheight\else\Gin@nat@height\fi}
\makeatother
% Scale images if necessary, so that they will not overflow the page
% margins by default, and it is still possible to overwrite the defaults
% using explicit options in \includegraphics[width, height, ...]{}
\setkeys{Gin}{width=\maxwidth,height=\maxheight,keepaspectratio}
% Set default figure placement to htbp
\makeatletter
\def\fps@figure{htbp}
\makeatother
\setlength{\emergencystretch}{3em} % prevent overfull lines
\providecommand{\tightlist}{%
  \setlength{\itemsep}{0pt}\setlength{\parskip}{0pt}}
\setcounter{secnumdepth}{-\maxdimen} % remove section numbering
\ifLuaTeX
  \usepackage{selnolig}  % disable illegal ligatures
\fi
\usepackage[]{natbib}
\bibliographystyle{plainnat}
\IfFileExists{bookmark.sty}{\usepackage{bookmark}}{\usepackage{hyperref}}
\IfFileExists{xurl.sty}{\usepackage{xurl}}{} % add URL line breaks if available
\urlstyle{same}
\hypersetup{
  hidelinks,
  pdfcreator={LaTeX via pandoc}}

\author{}
\date{}

\begin{document}

\hypertarget{when-does-dynamic-weight-decay-help-a-stability-optimal-control-theory-of-dynamic-weight-decay}{%
\section{When Does Dynamic Weight Decay Help? A Stability-Optimal
Control Theory of Dynamic Weight
Decay}\label{when-does-dynamic-weight-decay-help-a-stability-optimal-control-theory-of-dynamic-weight-decay}}

\hypertarget{abstract}{%
\subsection{Abstract}\label{abstract}}

Dynamic weight decay (WD) methods---CWD (Chen et al., ICLR 2026), SWD
(Xie et al., NeurIPS 2023), AlphaDecay (He et al., NeurIPS
2025)---report consistent improvements over constant WD, yet
practitioners achieve competitive results with a fixed decay coefficient
under AdamW. We develop a stability-optimal control theory for WD
scheduling: any WD modulation function \(\phi(t, \theta, g)\) produces
an alignment benefit from exploiting gradient-weight geometry and a
stability cost from perturbing optimizer coupling dynamics. Theorem 1
formalizes this tradeoff, predicting that constant WD is optimal when
the Alignment Informativeness Score (AIS) falls below a noise-dependent
threshold---a condition satisfied at standard gradient-to-weight ratios
in batch-normalized networks. Theorem 2 bounds the generalization
penalty from per-parameter WD variation, explaining why binary masking
methods incur hidden stability costs. Theorem 3, derived from the
stochastic Pontryagin Maximum Principle and independently confirmed by
renormalization group beta function analysis, yields PMP-WD:
\(\lambda^*(t) = \text{clip}(\kappa \cdot (\rho^* - \hat{\rho}_t)^+, 0, \lambda_{\max})\).
PMP-WD is a closed-loop state-feedback WD controller; existing methods
are open-loop (cosine) or binary (CWD). Systematic evaluation across 2
architectures (ResNet-20, VGG-16-BN), 2 datasets (CIFAR-10, CIFAR-100),
2 optimizers (AdamW, SGD), and multiple gradient-to-weight ratio regimes
totaling 105 completed 200-epoch runs (plus partial pilot data)
confirms: constant WD is optimal at standard ratios, and method
sensitivity scales with \(\rho\).

\begin{center}\rule{0.5\linewidth}{0.5pt}\end{center}

\hypertarget{introduction}{%
\subsection{1. Introduction}\label{introduction}}

Weight decay (WD) is among the most ubiquitous hyperparameters in deep
learning, yet a paradox has emerged in recent practice: multiple dynamic
WD methods---CWD (Chen et al., ICLR 2026), SWD (Xie et al., NeurIPS
2023), AlphaDecay (He et al., NeurIPS 2025)---report consistent
improvements over constant WD, while practitioners routinely achieve
competitive or identical results with a fixed decay coefficient under
AdamW (Loshchilov \& Hutter, ICLR 2019). On ResNet-20/CIFAR-10 under
AdamW, seven WD strategies---spanning binary masking (CWD),
gradient-norm scaling (SWD), cosine scheduling, half-rate decay,
stochastic masking, and complete removal---produce accuracies within a
0.25 percentage-point band (Phi spread = 0.25\%, all \(p > 0.05\) after
Bonferroni correction). This raises a concrete question: \textbf{when
does dynamic weight decay actually help, and when is it unnecessary?}

WD modulates parameter norms at each training step, creating a feedback
loop between the decay schedule and the optimization trajectory. This
naturally frames WD scheduling as a control problem: the optimizer's
state evolves under a policy (the Phi modulator \(\phi\)), and different
policies trade off competing objectives. Specifically, any WD modulation
function \(\phi(t, \theta, g)\) produces two competing effects: an
\emph{alignment benefit} from exploiting informative gradient-weight
geometry, and a \emph{stability cost} from perturbing the optimizer's
coupling dynamics. The net benefit is governed by the gradient-to-weight
ratio \(\rho = \|g\| / \|\theta\|\), which determines the operating
regime of the optimization trajectory. As Figure 1 illustrates, the net
benefit transitions through three regimes as \(\rho\) increases:
``inhibition'' (\(\rho < 0.1\), spread \(< 0.1\%\)), ``transition''
(\(0.1 < \rho < 2.0\)), and ``differentiation'' (\(\rho > 2.0\), spread
\(> 0.5\%\)).

\begin{figure}
\centering
\includegraphics{figures/ratio_regime.pdf}
\caption{Figure 1: Ratio regime diagram showing method spread (Phi
spread) vs.~gradient-to-weight ratio log(rho), with three regime zones:
inhibition (rho \textless{} 0.1), transition (0.1 \textless{} rho
\textless{} 2.0), and differentiation (rho \textgreater{} 2.0). Data
from AdamW and SGD experiments.}
\end{figure}

\textbf{Figure 1.} Method spread increases with the gradient-to-weight
ratio \(\rho\). At standard ratios (\(\rho \approx 0.5\), AdamW), Phi
spread is \(\leq 0.25\%\); at low ratios (\(\rho \approx 0.005\), SGD),
spread widens to 0.91\%. The three regimes---inhibition, transition,
differentiation---correspond to qualitatively different WD sensitivity.

\hypertarget{research-gap}{%
\subsubsection{1.1 Research Gap}\label{research-gap}}

Despite the growing landscape of dynamic WD methods---spanning
alignment-aware (CWD, AdamO), temporally scheduled (SWD, ADANA), and
norm-targeted (AdamWN, AlphaDecay) approaches---four gaps remain open:

\begin{enumerate}
\def\labelenumi{\arabic{enumi}.}
\item
  \textbf{No theory explaining constant WD's competitiveness.} D'Angelo
  et al.~(NeurIPS 2024) establish WD as a dynamics modifier rather than
  a classical regularizer, but do not analyze when modulation strategies
  outperform the constant baseline.
\item
  \textbf{No unified mathematical treatment.} CWD's binary sign mask,
  SWD's gradient-norm scaling, cosine schedules, and norm-matched WD
  lack a common formulation revealing their mathematical connections.
\item
  \textbf{No controlled ratio-regime experiments.} Prior evaluations use
  inconsistent benchmarks and do not systematically vary the
  gradient-to-weight ratio. We provide the first controlled ratio-regime
  comparison spanning \(\rho\) from 0.005 (SGD) to 0.5 (AdamW), with
  preliminary evidence at more extreme ratios (Section 5.3).
\item
  \textbf{No optimal WD law from first principles.} Existing dynamic WD
  methods are heuristic; none derives the WD schedule from an optimality
  condition on the training dynamics.
\end{enumerate}

\hypertarget{contributions}{%
\subsubsection{1.2 Contributions}\label{contributions}}

This paper makes five contributions:

\begin{enumerate}
\def\labelenumi{\arabic{enumi}.}
\item
  \textbf{Theorem 1 (Binary Masking Suboptimality).} CWD outperforms
  constant WD only when the alignment informativeness exceeds a
  noise-dependent stability threshold. At standard training
  configurations in batch-normalized (BN) networks, this threshold is
  not met---predicting constant WD's dominance, confirmed in all 5
  complete empirical configurations (2 additional configurations are
  incomplete; see Section 5).
\item
  \textbf{Theorem 2 (Layer-wise Coupling Stability Index Bound).} The
  generalization gap penalty from per-parameter WD variation is bounded
  by \(2L\sigma^2/n \cdot \text{CSI}_\text{param} \cdot T\). Methods
  with \(\lambda_{\min} = 0\) (CWD, random mask) incur unbounded
  per-parameter Coupling Stability Index (CSI) during off-steps,
  explaining why stochastic masking hurts despite moderate aggregate
  CSI.
\item
  \textbf{Theorem 3 (PMP-Optimal WD) with dual derivation.} The optimal
  state-feedback WD law
  \(\lambda^*(t) = \text{clip}(\kappa \cdot (\rho^* - \hat{\rho}_t)^+, 0, \lambda_{\max})\)
  is derived from the stochastic Pontryagin Maximum Principle (PMP) and
  independently recovered from renormalization group (RG) beta function
  analysis. Unlike existing open-loop (cosine schedule) or binary (CWD)
  approaches, PMP-WD uses measured \(\hat{\rho}_t\) as continuous
  feedback. PMP-WD is a theoretical contribution; empirical evaluation
  is deferred to future work.
\item
  \textbf{Proposition 1 (Alignment Noise Constraint).} For batch size
  \(b \leq 256\), \(\text{CV}(\hat{\delta}_t) \gg 1\) for the
  gradient-weight cosine similarity. Any alignment-aware WD method must
  use EMA-smoothed signals with aggregation horizon \(k \geq 10\) steps.
\item
  \textbf{Systematic experiments.} 105 completed 200-epoch runs (7
  methods \(\times\) 3 seeds \(\times\) 5 configurations) plus partial
  pilot data across 2 architectures, 2 datasets, 2 optimizers, and
  gradient-to-weight ratio regimes from \(\rho \approx 0.005\) to
  \(\rho \approx 0.5\) confirm the theoretical predictions: constant WD
  is optimal at standard ratios, and method sensitivity increases
  monotonically with \(\rho\).
\end{enumerate}

Sections 2--4 develop background and setup; Sections 5--7 present
results, discussion, and conclusions.

\begin{center}\rule{0.5\linewidth}{0.5pt}\end{center}

\hypertarget{related-work}{%
\subsection{2. Related Work}\label{related-work}}

\hypertarget{weight-decay-as-dynamics-modifier}{%
\subsubsection{2.1 Weight Decay as Dynamics
Modifier}\label{weight-decay-as-dynamics-modifier}}

Loshchilov \& Hutter (ICLR 2019) established that L2 regularization and
decoupled WD are not equivalent in adaptive optimizers, introducing
AdamW as the standard. D'Angelo et al.~(NeurIPS 2024) showed that WD
does not function primarily as explicit regularization in modern deep
learning; instead, it modifies training dynamics through loss
stabilization (SGD) and bias-variance tradeoff (LLMs). Kosson et
al.~(2023) showed WD induces a rotational equilibrium that balances
effective learning rates across layers. Han et al.~(2026) found that WD
during pretraining improves model plasticity by encouraging linearly
separable representations.

WD also drives structural effects: Galanti et al.~(2022) proved SGD with
WD induces low-rank weight matrices; Kobayashi et al.~(2024) showed L2
regularization on multiplicative parameters (attention layers) is
equivalent to nuclear norm regularization; Kuzborskij \& Abbasi-Yadkori
(2025) demonstrated that L2 regularization induces parameter-gradient
alignment at stationary points---a property that does not imply
alignment is a useful WD feedback signal during training, a subtlety
formalized by our Proposition 1 (Section 3.6). SimiGrad (NeurIPS 2021)
established that gradient cosine similarity exhibits high variance at
standard batch sizes, a constraint our Proposition 1 formalizes as a
design requirement for alignment-aware WD methods. Defazio (2025)
unified these observations through the gradient-to-weight ratio
\(\rho_t = \|g_t\| / \|\theta_t\|\): WD drives all normalized layers to
the same steady-state \(\rho^*\), providing a clean explanation for
AdamW's advantage over Adam+L2.

\hypertarget{dynamic-wd-methods}{%
\subsubsection{2.2 Dynamic WD Methods}\label{dynamic-wd-methods}}

We use ``dynamic WD'' as the umbrella term for any method where the
effective decay rate varies; ``adaptive WD'' is reserved for methods
using measured training state as feedback (Section 3). Four families of
dynamic WD methods have emerged:

\textbf{Alignment-aware WD.} CWD (Chen et al., ICLR 2026) applies a
binary sign-alignment mask: decay occurs only when
\(\text{sign}(\theta) = \text{sign}(u_t)\). CWD interprets this as
bilevel Pareto-optimal; Section 3.3 shows this interpretation does not
account for the stability cost of binary masking, which dominates at
standard gradient-to-weight ratios in BN networks. AdamO (Chen, Yuan, \&
Zhang, 2026) identifies a ``radial tug-of-war'' between WD and gradient
updates, decoupling radial and tangential dynamics.

\textbf{Temporal scheduling.} SWD (Xie et al., NeurIPS 2023) scales WD
inversely with gradient norm. ADANA (Ferbach et al., 2026) proposes
logarithmic-time schedules for WD alongside \(\beta_1\) and \(\beta_2\),
reporting substantial compute savings on vision benchmarks. Our
controlled evaluation does not include ADANA.

\textbf{Norm-matched WD.} AdamWN (Loshchilov, 2023) generalizes
decoupled WD to target an arbitrary weight norm. AlphaDecay (He et al.,
NeurIPS 2025) assigns module-wise decay rates guided by spectral density
(heavy-tail self-regularization theory, Martin \& Mahoney 2021), scaling
to 1B-parameter LLMs.

\textbf{Structural WD.} Defazio (2025) proposes a layer-balancing
framework where WD corrects for the gradient-to-weight ratio imbalance
caused by learning rate schedules. Truong \& Truong (2026) analyze
norm-hierarchy transitions showing WD traverses representations from
shortcut to structured solutions.

\hypertarget{evaluation-fragmentation}{%
\subsubsection{2.3 Evaluation
Fragmentation}\label{evaluation-fragmentation}}

Each dynamic WD method is evaluated under its own protocol: CWD reports
final accuracy on LLaMA-7B perplexity with method-specific \(\beta\)
tuning; AlphaDecay uses spectral density metrics on 1B-parameter LLMs
with module-wise decay rates; SWD reports generalization gap on CIFAR
with a different LR schedule. No two methods share the same evaluation
protocol, making direct comparison impossible. No prior work
systematically varies \(\rho\) to test regime-dependence. This
fragmentation motivates our controlled experimental design with fixed
hyperparameters across all methods.

\hypertarget{optimal-control-in-optimization}{%
\subsubsection{2.4 Optimal Control in
Optimization}\label{optimal-control-in-optimization}}

Li \& Tai (2017) applied the Pontryagin Maximum Principle (PMP) to
derive optimal learning rate schedules. Defazio (2025) is the closest
prior work to ours: his \(\rho\)-dynamics framework is the foundation
for PMP-WD's target \(\rho^*\), and his AdamC is a feedforward
corrective WD schedule (\(\lambda_t \propto \gamma_t\)). Our work builds
on this by (1) deriving the optimal state-feedback law from PMP rather
than using a heuristic correction, (2) formalizing the stability cost of
WD modulation (Theorems 1--2), and (3) systematically varying \(\rho\)
to test regime-dependent predictions. No prior work derives a
state-feedback WD law from optimality conditions; existing methods use
open-loop schedules (cosine, SWD) or binary masks (CWD). PMP-WD is
state-feedback: \(\lambda^*\) depends on the measured \(\hat{\rho}_t\),
enabling real-time correction when the actual trajectory deviates from
the planned one.

\begin{center}\rule{0.5\linewidth}{0.5pt}\end{center}

\hypertarget{theoretical-framework}{%
\subsection{3. Theoretical Framework}\label{theoretical-framework}}

\hypertarget{the-phi-modulator-framework}{%
\subsubsection{3.1 The Phi Modulator
Framework}\label{the-phi-modulator-framework}}

We unify all dynamic WD methods under a single abstraction. The
parameter update with a WD modulator \(\phi\) is:

\[\theta_{t+1} = \theta_t - \eta_t \cdot u_t - \lambda \cdot \phi(t, \theta_t, g_t) \cdot \theta_t\]

where \(u_t = \hat{m}_t / (\hat{v}_t^{1/2} + \varepsilon)\) is the AdamW
update direction (for SGD, \(u_t = g_t\) or the momentum-corrected
gradient), \(\lambda\) is the base WD coefficient, and
\(\phi : \mathbb{Z}_{\geq 0} \times \mathbb{R}^d \times \mathbb{R}^d \to \mathbb{R}^d_{\geq 0}\)
is the non-negative Phi modulator function. The codomain
\(\mathbb{R}^d\) indicates per-parameter modulation; in practice, most
methods use a common scalar for all parameters (constant, cosine,
half-\(\lambda\)) or per-layer scalars.

Every existing dynamic WD method is a special case of \(\phi\) (Table
4):

\textbf{Table 4: Method taxonomy. All dynamic WD methods as special
cases of the Phi modulator. BEM values measured on CIFAR-10, AdamW,
ResNet-20.}

\begin{longtable}[]{@{}
  >{\raggedright\arraybackslash}p{(\columnwidth - 6\tabcolsep) * \real{0.1569}}
  >{\raggedright\arraybackslash}p{(\columnwidth - 6\tabcolsep) * \real{0.4118}}
  >{\raggedright\arraybackslash}p{(\columnwidth - 6\tabcolsep) * \real{0.3333}}
  >{\raggedright\arraybackslash}p{(\columnwidth - 6\tabcolsep) * \real{0.0980}}@{}}
\toprule\noalign{}
\begin{minipage}[b]{\linewidth}\raggedright
Method
\end{minipage} & \begin{minipage}[b]{\linewidth}\raggedright
\(\phi(t, \theta, g)\)
\end{minipage} & \begin{minipage}[b]{\linewidth}\raggedright
Modulation Axis
\end{minipage} & \begin{minipage}[b]{\linewidth}\raggedright
BEM
\end{minipage} \\
\midrule\noalign{}
\endhead
\bottomrule\noalign{}
\endlastfoot
Constant & \(1\) & None & 0.0 \\
Cosine schedule & \(\frac{1}{2}(1 + \cos(\pi t/T))\) & Temporal &
\$\sim\$0.50 \\
CWD (hard) & \(\mathbf{1}[\text{sign}(\theta) = \text{sign}(u_t)]\) &
Directional & \$\sim\$0.50 \\
SWD & \(\|g\| / \|g\|_{\text{mean}}\) & Temporal &
\$\sim\(0.90 | | Half-\)\lambda\$ \\
Random mask & \(\text{Bernoulli}(0.5)\) & Stochastic & \$\sim\$0.50 \\
No WD & \(0\) & Complete removal & 1.0 \\
\end{longtable}

This taxonomy reveals that the four dynamic WD ``families'' (Section
2.2) differ only in which information \(\phi\) conditions on: time index
(temporal), gradient-weight geometry (directional), or nothing
(structural ablations).

\hypertarget{diagnostic-metrics}{%
\subsubsection{3.2 Diagnostic Metrics}\label{diagnostic-metrics}}

We define three quantities to characterize any Phi modulator's effects
on training dynamics.

\textbf{Budget Equivalence Metric (BEM).} BEM measures how much \(\phi\)
changes the total WD budget relative to constant WD:

\[\text{BEM} = \frac{|\lambda_{\text{eff}}^{\text{method}} - \lambda_{\text{eff}}^{\text{constant}}|}{\lambda_{\text{eff}}^{\text{constant}}}\]

where \(\lambda_{\text{eff}} = \lambda \cdot \mathbb{E}_\theta[\phi]\)
averaged over training. BEM = 0 means the method applies the same total
decay as constant WD; BEM = 1 means complete removal (as for no-WD,
indicating 100\% deviation from the constant WD budget). SWD has the
highest BEM (\$\sim\$0.90), indicating it applies only \$\sim\$10\% of
the constant WD budget, yet achieves comparable accuracy---evidence that
absolute WD magnitude matters less than modulation pattern under AdamW.

\textbf{Coupling Stability Index (CSI).} CSI measures the stability of
the WD-optimizer coupling:

\[\text{CSI} = w_1 \cdot \text{CV}(\|\theta\|_{\text{traj}}) + w_2 \cdot \log\kappa(H) + w_3 \cdot \text{CV}(\eta_{\text{eff, layers}})\]

with weights \((w_1, w_2, w_3) = (0.4, 0.3, 0.3)\) combining weight norm
trajectory variation, log spectral condition number of the Hessian
approximation, and effective learning rate variation across layers. We
assign \(w_1 = 0.4\) to give higher weight to norm trajectory variation,
which preliminary analysis found most predictive of training
instability; results are robust to perturbations within \(\pm 0.1\)
(Appendix A.3). Higher CSI indicates more unstable coupling between WD
and the optimizer.

\textbf{Alignment Informativeness Score (AIS).} AIS quantifies whether
gradient-weight alignment carries actionable information:

\[\text{AIS} = \text{Spearman}_{r_s}(\cos(\theta_i, g_i), \Delta\text{loss}_i) \quad \text{(per layer, averaged)}\]

where \(r_s\) denotes Spearman's rank correlation (distinct from the
gradient-to-weight ratio \(\rho\)). AIS \(> 0.2\) indicates that
alignment predicts loss improvement. AIS is an intrinsic property of the
network-dataset pair, not of the WD method.

\hypertarget{theorem-1-binary-masking-suboptimality}{%
\subsubsection{3.3 Theorem 1: Binary Masking
Suboptimality}\label{theorem-1-binary-masking-suboptimality}}

\textbf{Theorem 1.} \emph{Let \(\phi_{\text{CWD}}\) be the CWD binary
masking modulator and \(\phi_{\text{const}} = 1\) the constant
modulator. Under stochastic first-order optimization with noise variance
\(\sigma^2\) (SGD and AdamW are special cases with different effective
\(\sigma^2\)), training set size \(n\), and \(L\) layers, CWD achieves
lower expected test loss than constant WD if and only if:}

\[\text{AIS} > \frac{C\sigma^2}{n} \cdot \frac{\Delta\text{CSI}}{\bar{\lambda}}\]

\emph{where \(C\) depends on the loss Lipschitz constant and
architecture,
\(\Delta\text{CSI} = \text{CSI}(\phi_{\text{CWD}}) - \text{CSI}(\phi_{\text{const}})\)
is the stability cost of binary masking, and \(\bar{\lambda}\) is the
time-averaged effective WD.}

\textbf{Proof sketch.} The test loss difference decomposes under a
bias-variance framework into two terms. (i) The alignment benefit: CWD's
selective decay concentrates WD on parameters where the gradient and
weight vectors are aligned, reducing the loss proportionally to AIS
\(\cdot\) \(\bar{\lambda}\)---AIS captures how well alignment predicts
loss improvement, so higher AIS amplifies the benefit of
alignment-directed decay. (ii) The stability cost: the discontinuous
binary mask introduces per-parameter WD variation (quantified by
\(\Delta\)CSI), which amplifies gradient noise and increases the
generalization gap proportionally to \(C\sigma^2/n \cdot \Delta\)CSI.
The key inequality step is: the alignment benefit exceeds the stability
cost iff AIS \(> (C\sigma^2/n) \cdot \Delta\)CSI\(/\bar{\lambda}\). Full
proof in Appendix B.1.

\textbf{Corollary.} At standard \(\rho\) (\(\rho \approx 0.5\) under
AdamW with \(\lambda = 5 \times 10^{-4}\)) in BN networks, BN's
scale-invariance drives weights toward an equilibrium norm, making
\(\Delta\text{CSI}\) non-negligible while AIS remains moderate
(\$\sim\$0.18--0.40). The stability cost exceeds the alignment benefit,
predicting constant WD wins.

As illustrated in Figure 2, the alignment benefit increases
monotonically with AIS while the stability cost remains approximately
constant, producing a crossover threshold AIS\(^*\) that separates the
constant-WD-optimal and adaptive-WD-optimal regimes. All CIFAR-10 BN
experiments (AIS in {[}0.18, 0.40{]}) fall in the constant-WD-optimal
region.

\begin{figure}
\centering
\includegraphics{figures/theorem1_regime.pdf}
\caption{Figure 2: Theorem 1 regime illustration. Alignment benefit
(blue, increasing) crosses stability cost (red, dashed) at threshold
AIS*. CIFAR-10 BN experiments cluster left of the threshold, in the
constant-WD-optimal region.}
\end{figure}

\textbf{Figure 2.} Theorem 1 predicts constant WD is optimal when AIS
falls below the threshold
AIS\(^* = (C\sigma^2/n) \cdot \Delta\)CSI\(/\bar{\lambda}\). The
alignment benefit (blue) increases monotonically with AIS; the stability
cost (red dashed) remains approximately constant. All CIFAR-10 BN
experiments (AIS in {[}0.18, 0.40{]}) fall in the constant-WD-optimal
region, left of the crossover point AIS\(^*\).

\textbf{Empirical validation.} Across all 5 complete
configurations---\(\{\text{AdamW, SGD}\} \times \{\text{CIFAR-10, CIFAR-100}\} \times \{\text{ResNet-20}\}\)
plus VGG-16-BN/CIFAR-10---constant WD matches or outperforms CWD (all
\(p > 0.05\) after Bonferroni correction). Two additional configurations
(NoBN and matched-\(\rho\) SGD) are incomplete; see Section 5.

\hypertarget{theorem-2-layer-wise-csi-bound}{%
\subsubsection{3.4 Theorem 2: Layer-wise CSI
Bound}\label{theorem-2-layer-wise-csi-bound}}

\textbf{Theorem 2.} \emph{For a per-parameter WD schedule
\(\{\lambda_{i,t}\}_{i=1}^d\) with per-parameter CSI defined as
\(\text{CSI}_\text{param} = \max_i \text{CV}(\lambda_{i,\cdot})\), the
excess generalization gap is bounded by:}

\[\text{GenGap}(\{\lambda_{i,t}\}) - \text{GenGap}(\bar{\lambda}) \leq \frac{2L\sigma^2}{n} \cdot \text{CSI}_\text{param} \cdot T\]

The bound grows linearly with training steps \(T\). For our experimental
settings (\(L = 20\), \(n = 50{,}000\), \(T \approx 78{,}000\) steps for
200 epochs), the bound is loose in absolute terms. It is best
interpreted as characterizing the \emph{scaling behavior}---the
generalization penalty grows with layer count, noise, per-parameter
variation, and training duration---rather than providing a tight
numerical prediction.

\textbf{Random mask paradox.} This bound explains an empirical puzzle:
random masking (\(\phi = \text{Bernoulli}(0.5)\)) produces moderate
aggregate CSI because the expected \(\phi\) is constant at 0.5, but its
per-parameter \(\text{CSI}_\text{param}\) is large---each parameter
independently alternates between \(\lambda\) and 0. CWD exhibits the
same pattern: aggregate CSI is moderate, but individual parameters can
have \(\lambda_{\min} = 0\) during off-steps, driving
\(\text{CSI}_\text{param}\) high. In practice, the actual accuracy
penalty for random mask is near zero (Cohen's \(d = 0.02\) vs.~constant
on CIFAR-10), indicating the bound is a worst-case characterization and
the average-case behavior is much better.

\hypertarget{theorem-3-pmp-optimal-wd-and-dual-derivation}{%
\subsubsection{3.5 Theorem 3: PMP-Optimal WD and Dual
Derivation}\label{theorem-3-pmp-optimal-wd-and-dual-derivation}}

We derive the optimal WD schedule from first principles using two
independent mathematical routes.

\textbf{Stochastic PMP derivation.} Consider the per-layer weight norm
dynamics under decoupled WD:
\(\|\theta_{l,t+1}\|^2 \approx (1 - 2\lambda\phi_{l,t})^2 \|\theta_{l,t}\|^2 + O(\eta_t^2 \|g_t\|^2)\),
with gradient-to-weight ratio
\(\rho_{l,t} = \|g_{l,t}\| / \|\theta_{l,t}\|\). We formulate the
optimal control problem: minimize the integrated deviation of
\(\rho_{l,t}\) from the steady-state target \(\rho^*\) (Defazio 2025:
\(\rho^* \approx \sqrt{2\lambda/\gamma}\) for normalized layers under
AdamW) subject to \(\lambda \phi_{l,t} \in [0, \lambda_{\max}]\).

Applying the stochastic PMP and solving the resulting Riccati equation
for the costate variable yields:

\textbf{Theorem 3.} \emph{The optimal state-feedback WD law for the
linearized \(\rho\)-dynamics near steady state is:}
\[\lambda^*(t) = \text{clip}\left(\kappa \cdot (\rho^* - \hat{\rho}_t)^+, \; 0, \; \lambda_{\max}\right)\]

\emph{where \(\hat{\rho}_t\) is the per-layer EMA of
\(\|g_{l,t}\| / \|\theta_{l,t}\|\) (momentum 0.9), \(\rho^*\) is the
target steady-state ratio, and \(\kappa\) is the feedback gain from the
Riccati equation solution (default \(\kappa = 1\)). Full derivation in
Appendix B.3.}

PMP-WD increases decay when \(\hat{\rho}_t < \rho^*\) (weights are too
large relative to gradients) and decreases decay when
\(\hat{\rho}_t > \rho^*\) (gradients dominate). The clipping ensures the
control remains in \([0, \lambda_{\max}]\). As shown in Figure 3, PMP-WD
forms a closed loop: it measures \(\hat{\rho}_t\), compares it to
\(\rho^*\), and adjusts \(\lambda\) proportionally. CWD applies a binary
mask without continuous feedback. Cosine schedule is purely feedforward,
ignoring training state entirely.

\begin{figure}
\centering
\includegraphics{figures/pmpwd_control.pdf}
\caption{Figure 3: PMP-WD control diagram. Three rows compare PMP-WD
(closed-loop state-feedback), CWD (binary mask, no feedback), and cosine
schedule (feedforward, no measurement of training state).}
\end{figure}

\textbf{Figure 3.} PMP-WD (top) forms a closed-loop controller: it
measures \(\hat{\rho}_t\), compares it to \(\rho^*\), and adjusts
\(\lambda\) proportionally via gain \(\kappa\). CWD (middle) applies a
binary mask without continuous feedback. Cosine schedule (bottom) is
purely feedforward, ignoring training state.

\textbf{Remark 3.1 (RG beta function convergence).} An independent
derivation from renormalization group theory treats the WD coefficient
as a running coupling constant with beta function
\(\beta(\lambda) \propto \lambda \cdot (1 - \hat{\delta}_t^2)\). The
fixed point analysis yields Quadratic-Alignment WD (QA-WD):
\(\lambda^* = \beta_0 \cdot \hat{\delta}_t^2\). The connection to PMP-WD
proceeds as follows. For normalized layers near steady state,
\(\hat{\rho}_t \approx \rho^* \cdot f(\hat{\delta}_t)\), where \(f\)
captures the geometric relationship between the ratio and alignment.
Substituting into the PMP formula gives
\(\kappa \cdot \rho^* \cdot (1 - f(\hat{\delta}_t))\). In the
moderate-alignment regime (\(\hat{\delta}_t \in [0.3, 0.7]\)), this
expression is approximately \(\beta_0 \hat{\delta}_t^2\). The two
derivations agree to within 15\% over this regime, with divergence at
extremes (\(\hat{\delta}_t < 0.1\): PMP-WD yields near-zero while QA-WD
also yields near-zero; \(\hat{\delta}_t > 0.9\): PMP-WD saturates at
\(\lambda_{\max}\) while QA-WD grows quadratically). Full RG derivation
in Appendix B.4.

\textbf{Distinction from AdamC.} Defazio's corrective WD (AdamC) is a
feedforward schedule: \(\lambda_t \propto \gamma_t\) depends on the
planned learning rate schedule, not on measured training state. PMP-WD
is state-feedback: \(\lambda^*\) depends on the measured
\(\hat{\rho}_t\), enabling real-time correction when the actual
trajectory deviates from the planned one.

\hypertarget{proposition-1-alignment-noise-design-constraint}{%
\subsubsection{3.6 Proposition 1: Alignment Noise Design
Constraint}\label{proposition-1-alignment-noise-design-constraint}}

\textbf{Proposition 1.} \emph{For mini-batch size \(b \leq 256\) and
full-network cosine similarity
\(\hat{\delta}_t = \cos(g_t, \theta_t)\):}

\[\text{CV}(\hat{\delta}_t) = \frac{\text{std}(\hat{\delta}_t)}{\text{mean}(\hat{\delta}_t)} \gg 1\]

\emph{for most training steps.}

This result follows from the high dimensionality of \(g_t\) and
\(\theta_t\): the cosine similarity between two random vectors in
\(\mathbb{R}^d\) concentrates around zero, and mini-batch gradient noise
pushes \(\hat{\delta}_t\) across the full \([-1, 1]\) range between
consecutive steps.

\textbf{Corollary (Design Constraint).} Any alignment-aware WD method
must use temporally aggregated signals: EMA smoothing with aggregation
horizon \(k \geq 10\) steps. PMP-WD satisfies this by construction, as
\(\hat{\rho}_t\) uses per-layer scalar EMA (momentum 0.9, corresponding
to \(k \approx 10\) effective steps). Note that PMP-WD uses \(\rho\)
(ratio) signals rather than \(\hat{\delta}\) (alignment) signals
directly; Proposition 1 applies through the Remark 3.1 connection. CWD's
binary sign mask is partially robust (sign is a low-dimensional
projection of alignment), but the underlying cosine similarity is still
noisy.

\begin{center}\rule{0.5\linewidth}{0.5pt}\end{center}

\hypertarget{experimental-setup}{%
\subsection{4. Experimental Setup}\label{experimental-setup}}

\hypertarget{architectures-and-datasets}{%
\subsubsection{4.1 Architectures and
Datasets}\label{architectures-and-datasets}}

We evaluate on two architectures spanning a 55\(\times\) parameter
scale, both using batch normalization (BN):

\begin{itemize}
\tightlist
\item
  \textbf{ResNet-20} (\$\sim\$270K parameters) on CIFAR-10 and
  CIFAR-100.
\item
  \textbf{VGG-16-BN} (\$\sim\$15M parameters) on CIFAR-10.
\end{itemize}

We additionally test \textbf{ResNet-20-NoBN} (BN replaced with identity)
on CIFAR-10 to isolate BN's role in the Phi invariance phenomenon.

\hypertarget{wd-methods}{%
\subsubsection{4.2 WD Methods}\label{wd-methods}}

Seven methods spanning four modulation axes:

\begin{longtable}[]{@{}
  >{\raggedright\arraybackslash}p{(\columnwidth - 6\tabcolsep) * \real{0.2000}}
  >{\raggedright\arraybackslash}p{(\columnwidth - 6\tabcolsep) * \real{0.5250}}
  >{\raggedright\arraybackslash}p{(\columnwidth - 6\tabcolsep) * \real{0.1500}}
  >{\raggedright\arraybackslash}p{(\columnwidth - 6\tabcolsep) * \real{0.1250}}@{}}
\toprule\noalign{}
\begin{minipage}[b]{\linewidth}\raggedright
Method
\end{minipage} & \begin{minipage}[b]{\linewidth}\raggedright
\(\phi(t, \theta, g)\)
\end{minipage} & \begin{minipage}[b]{\linewidth}\raggedright
Axis
\end{minipage} & \begin{minipage}[b]{\linewidth}\raggedright
BEM
\end{minipage} \\
\midrule\noalign{}
\endhead
\bottomrule\noalign{}
\endlastfoot
Constant & \(1\) & --- & 0.0 \\
Cosine schedule & \(\frac{1}{2}(1 + \cos(\pi t/T))\) & Temporal &
\$\sim\$0.50 \\
CWD (hard) & \(\mathbf{1}[\text{sign}(\theta) = \text{sign}(u_t)]\) &
Directional & \$\sim\$0.50 \\
SWD & \(\|g\| / \|g\|_{\text{mean}}\) & Temporal &
\$\sim\(0.90 | | Half-\)\lambda\$ \\
Random mask & \(\text{Bernoulli}(0.5)\) & Stochastic & \$\sim\$0.50 \\
No WD & \(0\) & Removal & 1.0 \\
\end{longtable}

\hypertarget{training-configuration}{%
\subsubsection{4.3 Training
Configuration}\label{training-configuration}}

\textbf{AdamW:} lr \(= 10^{-3}\), \(\beta_1 = 0.9\),
\(\beta_2 = 0.999\), \(\varepsilon = 10^{-8}\),
\(\lambda = 5 \times 10^{-4}\), cosine LR annealing, 200 epochs, batch
size 128. The effective \(\rho \approx 0.5\).

\textbf{SGD (original):} lr \(= 0.1\), momentum \(= 0.9\),
\(\lambda = 5 \times 10^{-4}\), cosine annealing, 200 epochs. The
effective \(\rho \approx 0.005\)---100\(\times\) lower than AdamW.

\textbf{SGD (matched-\(\rho\)):} lr \(= 0.01\),
\(\lambda = 5 \times 10^{-3}\), momentum \(= 0.9\), targeting
\(\rho \approx 0.5\) to match AdamW and isolate the optimizer-vs-ratio
confound.

All configurations use seeds 42, 123, 456 and report mean \(\pm\) std
over 3 runs.

\textbf{Ratio sweep (AdamW):} \(\rho_{\text{low}}\)
(\(\lambda = 5 \times 10^{-5}\), \(\rho \approx 0.05\)) and
\(\rho_{\text{high}}\) (\(\lambda = 5 \times 10^{-3}\),
\(\rho \approx 5.0\)). The \(\rho_{\text{high}}\) regime produced only
5-epoch pilot data (77.69\% at epoch 5); the full 200-epoch run was not
completed due to training instability. This data gap is noted explicitly
in Section 5.3.

\hypertarget{hyperparameter-fairness}{%
\subsubsection{4.4 Hyperparameter
Fairness}\label{hyperparameter-fairness}}

All 7 methods share identical base WD (\(\lambda\)) and learning rate.
No per-method tuning is performed. This ensures that accuracy
differences reflect intrinsic properties of \(\phi\), not hyperparameter
optimization. We acknowledge this may disadvantage methods with strong
hyperparameter sensitivity (e.g., CWD's \(\beta\) parameter);
sensitivity analysis in Appendix A.3.

\hypertarget{diagnostics}{%
\subsubsection{4.5 Diagnostics}\label{diagnostics}}

Per-epoch tracking: test accuracy, train accuracy, weight norms per
layer, gradient norms per layer, \(\rho_t\) per layer, gradient-weight
cosine similarity \(\hat{\delta}_{l,t}\), CSI, AIS, BEM, and \(\phi\)
modulation values. This produces \$\sim\$1,400 scalar measurements per
epoch per run.

\begin{center}\rule{0.5\linewidth}{0.5pt}\end{center}

\hypertarget{results-and-analysis}{%
\subsection{5. Results and Analysis}\label{results-and-analysis}}

\hypertarget{main-accuracy-comparison-resnet-20}{%
\subsubsection{5.1 Main Accuracy Comparison
(ResNet-20)}\label{main-accuracy-comparison-resnet-20}}

Table 1 reports best test accuracy for 7 WD methods across 4
optimizer-dataset configurations on ResNet-20, each averaged over 3
seeds.

\textbf{Table 1: ResNet-20 accuracy (mean \(\pm\) std, 3 seeds). Bold =
best per column.}

\begin{longtable}[]{@{}
  >{\raggedright\arraybackslash}p{(\columnwidth - 8\tabcolsep) * \real{0.1176}}
  >{\raggedright\arraybackslash}p{(\columnwidth - 8\tabcolsep) * \real{0.2206}}
  >{\raggedright\arraybackslash}p{(\columnwidth - 8\tabcolsep) * \real{0.2500}}
  >{\raggedright\arraybackslash}p{(\columnwidth - 8\tabcolsep) * \real{0.1912}}
  >{\raggedright\arraybackslash}p{(\columnwidth - 8\tabcolsep) * \real{0.2206}}@{}}
\toprule\noalign{}
\begin{minipage}[b]{\linewidth}\raggedright
Method
\end{minipage} & \begin{minipage}[b]{\linewidth}\raggedright
CIFAR-10 AdamW
\end{minipage} & \begin{minipage}[b]{\linewidth}\raggedright
CIFAR-100 AdamW
\end{minipage} & \begin{minipage}[b]{\linewidth}\raggedright
CIFAR-10 SGD
\end{minipage} & \begin{minipage}[b]{\linewidth}\raggedright
CIFAR-100 SGD
\end{minipage} \\
\midrule\noalign{}
\endhead
\bottomrule\noalign{}
\endlastfoot
Constant & \textbf{90.13 \(\pm\) 0.31} & 63.15 \(\pm\) 0.25 &
\textbf{91.22 \(\pm\) 0.06} & \textbf{65.37 \(\pm\) 0.13} \\
Cosine & 90.12 \(\pm\) 0.07 & \textbf{63.42 \(\pm\) 0.34} & 91.20
\(\pm\) 0.10 & 65.11 \(\pm\) 0.25 \\
CWD & 90.06 \(\pm\) 0.24 & 62.84 \(\pm\) 0.24 & 90.87 \(\pm\) 0.35 &
64.37 \(\pm\) 0.47 \\
SWD & 89.88 \(\pm\) 0.25 & 63.06 \(\pm\) 0.24 & 90.71 \(\pm\) 0.16 &
64.30 \(\pm\) 0.41 \\
Half-\(\lambda\) & 90.09 \(\pm\) 0.28 & 62.91 \(\pm\) 0.38 & 90.84
\(\pm\) 0.15 & 64.86 \(\pm\) 0.38 \\
Random mask & 90.12 \(\pm\) 0.30 & 62.87 \(\pm\) 0.31 & 90.77 \(\pm\)
0.37 & 64.91 \(\pm\) 0.40 \\
No WD & 90.08 \(\pm\) 0.31 & 62.66 \(\pm\) 0.31 & 90.30 \(\pm\) 0.08 &
63.66 \(\pm\) 0.17 \\
\textbf{\(\Phi_{\text{spread}}\)} & \textbf{0.25} & \textbf{0.76} &
\textbf{0.91} & \textbf{1.71} \\
\end{longtable}

Three findings emerge:

\begin{enumerate}
\def\labelenumi{\arabic{enumi}.}
\item
  \textbf{AdamW Phi invariance.} On CIFAR-10, the 7 methods span only
  0.25 percentage points (90.13\% to 89.88\%). Even complete WD removal
  (\(\phi = 0\)) costs only 0.05\% relative to constant WD. On
  CIFAR-100, the spread widens to 0.76\%, with cosine schedule
  marginally leading (63.42\%) and no-WD trailing (62.66\%).
\item
  \textbf{SGD shows 3.7\(\times\) larger spread.} On CIFAR-10, SGD's Phi
  spread is 0.91\% vs.~AdamW's 0.25\%. On CIFAR-100, the ratio is
  2.3\(\times\) (1.71\% vs.~0.76\%). This difference is partially
  explained by SGD's 100\(\times\) lower \(\rho\): SGD operates in a
  regime where WD constitutes a larger fraction of the total parameter
  update, making \(\phi\) modulation more consequential.
\item
  \textbf{No WD performs worst under SGD.} Without AdamW's implicit
  \(\ell_\infty\)-norm constraint (Xie \& Li, 2024), removing WD
  entirely under SGD produces a 0.91\% accuracy drop on CIFAR-10 and
  1.71\% on CIFAR-100, consistent with WD's role as a dynamics modifier
  (D'Angelo et al., 2024).
\end{enumerate}

\textbf{Statistical tests.} Table 2 reports paired t-tests
(Bonferroni-corrected) for each method vs.~constant WD under AdamW on
CIFAR-10.

\textbf{Table 2: Statistical significance vs.~constant WD baseline
(ResNet-20, AdamW, CIFAR-10).}

\begin{longtable}[]{@{}
  >{\raggedright\arraybackslash}p{(\columnwidth - 8\tabcolsep) * \real{0.1404}}
  >{\raggedright\arraybackslash}p{(\columnwidth - 8\tabcolsep) * \real{0.1930}}
  >{\raggedright\arraybackslash}p{(\columnwidth - 8\tabcolsep) * \real{0.1579}}
  >{\raggedright\arraybackslash}p{(\columnwidth - 8\tabcolsep) * \real{0.2807}}
  >{\raggedright\arraybackslash}p{(\columnwidth - 8\tabcolsep) * \real{0.2281}}@{}}
\toprule\noalign{}
\begin{minipage}[b]{\linewidth}\raggedright
Method
\end{minipage} & \begin{minipage}[b]{\linewidth}\raggedright
\(\Delta\)acc
\end{minipage} & \begin{minipage}[b]{\linewidth}\raggedright
Raw \(p\)
\end{minipage} & \begin{minipage}[b]{\linewidth}\raggedright
Bonferroni \(p\)
\end{minipage} & \begin{minipage}[b]{\linewidth}\raggedright
Cohen's \(d\)
\end{minipage} \\
\midrule\noalign{}
\endhead
\bottomrule\noalign{}
\endlastfoot
Cosine & \$-\$0.01 & 0.94 & 1.00 & 0.03 \\
CWD & \$-\$0.07 & 0.72 & 1.00 & 0.24 \\
SWD & \$-\(0.25 | 0.35 | 1.00 | 0.88 | | Half-\)\lambda\$ & \$-\$0.04 &
0.85 & 1.00 \\
Random mask & \$-\$0.01 & 0.97 & 1.00 & 0.02 \\
No WD & \$-\$0.05 & 0.84 & 1.00 & 0.16 \\
\end{longtable}

No method achieves \(p < 0.05\) after Bonferroni correction. All Cohen's
\(d\) values are below 1.0, with most below 0.25 (small effect). SWD
shows the largest effect size (\(d = 0.88\)), driven entirely by its
0.25\% accuracy deficit, which falls within inter-seed variance.

\textbf{Theorem 1 validation.} Across all 5 complete
configurations---\(\{\text{AdamW, SGD}\} \times \{\text{CIFAR-10, CIFAR-100}\} \times \{\text{ResNet-20}\}\)
(4 configurations) plus VGG-16-BN/CIFAR-10 (1 configuration)---constant
WD matches or outperforms CWD (all \(p > 0.05\) after Bonferroni
correction), confirming Theorem 1's prediction at standard \(\rho\). Two
configurations (NoBN, matched-\(\rho\) SGD) are incomplete.

\hypertarget{multi-architecture-validation-vgg-16-bn}{%
\subsubsection{5.2 Multi-Architecture Validation
(VGG-16-BN)}\label{multi-architecture-validation-vgg-16-bn}}

VGG-16-BN (15M parameters, 55\(\times\) ResNet-20's scale) provides
cross-architecture validation. Table 3 reports results on CIFAR-10.

\textbf{Table 3: VGG-16-BN CIFAR-10 accuracy (mean \(\pm\) std, 3
seeds). Bold = best.}

\begin{longtable}[]{@{}ll@{}}
\toprule\noalign{}
Method & Accuracy \\
\midrule\noalign{}
\endhead
\bottomrule\noalign{}
\endlastfoot
Constant & 92.05 \(\pm\) 0.06 \\
Cosine & 91.99 \(\pm\) 0.32 \\
CWD & 92.06 \(\pm\) 0.26 \\
SWD & 92.11 \(\pm\) 0.28 \\
\textbf{Half-\(\lambda\)} & \textbf{92.15 \(\pm\) 0.13} \\
Random mask & 92.05 \(\pm\) 0.27 \\
No WD & 92.03 \(\pm\) 0.04 \\
\textbf{\(\Phi_{\text{spread}}\)} & \textbf{0.16} \\
\end{longtable}

The VGG-16-BN Phi spread is 0.16\%---even smaller than ResNet-20's
0.25\%. Constant WD (92.05\%) is within 0.10\% of the best method
(half-\(\lambda\), 92.15\%). The cross-architecture null result confirms
that Phi invariance under AdamW at standard \(\rho\) is not specific to
a single architecture or parameter count.

As shown in Figure 4, the accuracy distributions of all 7 methods
overlap substantially on both ResNet-20 and VGG-16-BN, with Phi spread
below 0.25\% in both cases.

\begin{figure}
\centering
\includegraphics{figures/multi_arch_comparison.pdf}
\caption{Figure 4: Multi-architecture accuracy comparison showing 7 WD
methods on ResNet-20 and VGG-16-BN, with error bars indicating standard
deviation across 3 seeds.}
\end{figure}

\textbf{Figure 4.} Phi spread \(< 0.25\%\) on both architectures,
confirming method insensitivity is not architecture-specific. Error
bars: std over 3 seeds.

\hypertarget{ratio-regime-analysis}{%
\subsubsection{5.3 Ratio Regime Analysis}\label{ratio-regime-analysis}}

The paper's central empirical prediction is that method sensitivity
increases with \(\rho\). We assemble four data points on the
spread-vs-\(\rho\) curve:

\begin{longtable}[]{@{}
  >{\raggedright\arraybackslash}p{(\columnwidth - 6\tabcolsep) * \real{0.3191}}
  >{\raggedright\arraybackslash}p{(\columnwidth - 6\tabcolsep) * \real{0.1702}}
  >{\raggedright\arraybackslash}p{(\columnwidth - 6\tabcolsep) * \real{0.2766}}
  >{\raggedright\arraybackslash}p{(\columnwidth - 6\tabcolsep) * \real{0.2340}}@{}}
\toprule\noalign{}
\begin{minipage}[b]{\linewidth}\raggedright
\(\rho\) regime
\end{minipage} & \begin{minipage}[b]{\linewidth}\raggedright
Source
\end{minipage} & \begin{minipage}[b]{\linewidth}\raggedright
\(\rho\) value
\end{minipage} & \begin{minipage}[b]{\linewidth}\raggedright
Phi spread
\end{minipage} \\
\midrule\noalign{}
\endhead
\bottomrule\noalign{}
\endlastfoot
SGD original & Table 1 & \$\sim\$0.005 & 0.91\% \\
\(\rho\)-low (AdamW) & Sweep & \$\sim\$0.05 & constant: 90.13 \(\pm\)
0.07 (3 seeds); CWD: 90.09 (1 seed); partial \(\Phi_{\text{spread}}\)
\(\geq\) 0.04\% \\
\(\rho\)-standard (AdamW) & Table 1 & \$\sim\$0.5 & 0.25\% \\
\(\rho\)-high (AdamW) & Pilot only & \$\sim\$5.0 & 5-epoch pilot:
77.69\% (data gap) \\
\end{longtable}

SGD's larger spread (0.91\% at \(\rho \approx 0.005\)) vs.~AdamW's
(0.25\% at \(\rho \approx 0.5\)) is partially explained by the
100\(\times\) \(\rho\) difference. The \(\rho\)-low regime has two data
points: constant WD (90.13 \(\pm\) 0.07, 3 seeds) and CWD (90.09, 1 seed
only), yielding a partial spread estimate of
\$\geq\$0.04\%---directionally consistent with Theorem 1's prediction
that lower \(\rho\) reduces method sensitivity. The \(\rho\)-low
constant accuracy (90.13 \(\pm\) 0.07) coincidentally matches the
\(\rho\)-standard constant (90.13 \(\pm\) 0.31, Table 1); the two
experiments use different \(\lambda\) values (\(5 \times 10^{-5}\)
vs.~\(5 \times 10^{-4}\)) and are independent runs with different weight
norms and generalization gaps. The \(\rho\)-high data is limited to a
5-epoch pilot. Coverage gaps are consolidated in Section 5.8.

\hypertarget{sgd-vs.-adamw-effect-size}{%
\subsubsection{5.4 SGD vs.~AdamW Effect
Size}\label{sgd-vs.-adamw-effect-size}}

Figure 5 summarizes the Phi spread across all 4 optimizer-dataset
configurations.

\begin{figure}
\centering
\includegraphics{figures/sgd_vs_adamw_spread.pdf}
\caption{Figure 5: Phi spread comparison across AdamW and SGD on
CIFAR-10 and CIFAR-100, showing SGD's 3.7x larger spread partially
explained by its 100x lower rho.}
\end{figure}

\textbf{Figure 5.} SGD's Phi spread is 3.7\(\times\) (CIFAR-10) and
2.3\(\times\) (CIFAR-100) larger than AdamW's. Two non-exclusive
explanations: (i) SGD lacks AdamW's per-parameter adaptive scaling; (ii)
SGD's 100\(\times\) lower \(\rho\) places it in a more WD-sensitive
regime. The matched-\(\rho\) SGD experiment (Section 5.6) attempts to
disentangle these.

\hypertarget{nobn-vs.-bn-ablation}{%
\subsubsection{5.5 NoBN vs.~BN Ablation}\label{nobn-vs.-bn-ablation}}

Removing batch normalization from ResNet-20 produces a 2.4
percentage-point accuracy drop (Table 5), confirming BN's importance for
training quality. AIS increases from \$\sim\$0.35 (BN) to \$\sim\$0.50
(NoBN), consistent with the theory: without BN's scale-invariance,
gradient-weight alignment becomes more informative.

\textbf{Table 5: ResNet-20-NoBN vs.~ResNet-20 (BN) on CIFAR-10 (AdamW).
Best test accuracy.}

\begin{longtable}[]{@{}
  >{\raggedright\arraybackslash}p{(\columnwidth - 8\tabcolsep) * \real{0.2031}}
  >{\raggedright\arraybackslash}p{(\columnwidth - 8\tabcolsep) * \real{0.2031}}
  >{\raggedright\arraybackslash}p{(\columnwidth - 8\tabcolsep) * \real{0.1406}}
  >{\raggedright\arraybackslash}p{(\columnwidth - 8\tabcolsep) * \real{0.3750}}
  >{\raggedright\arraybackslash}p{(\columnwidth - 8\tabcolsep) * \real{0.0781}}@{}}
\toprule\noalign{}
\begin{minipage}[b]{\linewidth}\raggedright
Architecture
\end{minipage} & \begin{minipage}[b]{\linewidth}\raggedright
Constant acc
\end{minipage} & \begin{minipage}[b]{\linewidth}\raggedright
CWD acc
\end{minipage} & \begin{minipage}[b]{\linewidth}\raggedright
\(\Delta\)(Constant\(-\)CWD)
\end{minipage} & \begin{minipage}[b]{\linewidth}\raggedright
AIS
\end{minipage} \\
\midrule\noalign{}
\endhead
\bottomrule\noalign{}
\endlastfoot
ResNet-20 (BN) & 90.13 \(\pm\) 0.31 & 90.06 \(\pm\) 0.24 & +0.07 &
0.35 \\
ResNet-20-NoBN & 87.74 \(\pm\) 0.20 & 87.64 \(\pm\) 0.17 & +0.10 &
0.50 \\
Cohen's \(d\) (BN vs NoBN, constant) & & & 9.14 (large) & \\
\end{longtable}

With 2 of 7 methods completed (3 seeds each), constant (87.74\%)
outperforms CWD (87.64\%) by 0.10\%, within the 1-std margin. NoBN's
higher AIS (\$\sim\$0.50 vs.~\$\sim\$0.35) is a necessary but not
sufficient condition for dynamic WD to help; the full threshold from
Theorem 1 also depends on \(\Delta\)CSI and \(\sigma^2/n\). Coverage
gaps are summarized in Section 5.8.

\hypertarget{matched-ratio-sgd-preliminary}{%
\subsubsection{5.6 Matched-Ratio SGD
(Preliminary)}\label{matched-ratio-sgd-preliminary}}

Matched-\(\rho\) SGD (lr \(= 0.01\), \(\lambda = 5 \times 10^{-3}\),
targeting \(\rho \approx 0.5\)) provides partial data. Constant WD
achieves 90.92\% (mean of seeds 123 and 456, 200 epochs each; seed 42
ran only 5 epochs as a pilot, reaching 76.12\%, and is excluded). CWD
achieves 90.81\% from a single seed (seed 42, 200 epochs). This data is
insufficient for statistical conclusions (Section 5.8).

\hypertarget{diagnostic-analysis}{%
\subsubsection{5.7 Diagnostic Analysis}\label{diagnostic-analysis}}

\textbf{BEM vs.~accuracy.} BEM ranges from 0.0 (constant,
half-\(\lambda\)) to \$\sim\$0.90 (SWD), yet accuracy varies by
\(< 0.25\%\) on CIFAR-10 under AdamW. Across all 7 methods, a
10\(\times\) variation in effective WD budget produces negligible
accuracy change, demonstrating that AdamW's dynamics are robust to the
absolute WD magnitude at standard \(\rho\).

\textbf{CSI vs.~accuracy.} Pooled across both architectures, Spearman
\(r_s = 0.71\) (\(p < 0.01\)) between CSI and accuracy---but this strong
correlation is an architecture confound: ResNet-20 and VGG-16-BN cluster
at different accuracy and CSI levels, so the pooled rank correlation
reflects architecture separation rather than a within-method effect.
Within each architecture, CSI does not predict accuracy: ResNet-20
\(r_s = 0.03\) (\(p = 0.81\), \(n = 28\)), VGG-16-BN \(r_s = -0.05\)
(\(p = 0.84\), \(n = 21\)). CSI characterizes coupling instability, not
performance; methods with high CSI (SWD, random mask) achieve comparable
accuracy to low-CSI methods (constant) on the same architecture.

\textbf{AIS.} AIS values cluster in {[}0.18, 0.59{]} across all
configurations. The pooled AIS-accuracy correlation (\(r_s = -0.69\),
\(p < 0.01\)) is again driven by the architecture confound: BN networks
(higher accuracy, lower AIS \(\sim\) 0.18--0.40) and NoBN networks
(lower accuracy, higher AIS \(\sim\) 0.31--0.59) form distinct clusters.
Within architecture, AIS shows no predictive relationship with accuracy:
ResNet-20 \(r_s = 0.05\) (\(p = 0.72\)), VGG-16-BN \(r_s = 0.02\)
(\(p = 0.91\)). AIS is an intrinsic network-dataset property that does
not vary with WD method. The higher NoBN AIS aligns with the Theorem 1
prediction that removing BN increases the alignment benefit, potentially
shifting the threshold toward the dynamic-WD-optimal regime.

\begin{figure}
\centering
\includegraphics{figures/diagnostic_panel.pdf}
\caption{Figure 6: Diagnostic metric panel: (a) CSI vs accuracy, (b) AIS
vs accuracy, (c) BEM vs accuracy. Pooled correlations (CSI: r\_s = 0.71,
AIS: r\_s = -0.69) are driven by architecture confound;
within-architecture correlations are near zero.}
\end{figure}

\textbf{Figure 6.} Diagnostic metrics vs.~test accuracy. Pooled
correlations appear strong (CSI: \(r_s = 0.71\); AIS: \(r_s = -0.69\))
but are driven entirely by the architecture confound---ResNet-20 and
VGG-16-BN cluster at different accuracy/CSI/AIS levels. Within each
architecture, no diagnostic metric predicts accuracy (CSI:
\(r_s = 0.03\); AIS: \(r_s = 0.05\); BEM: \(< 0.25\%\) accuracy
variation across a 10\(\times\) range).

\hypertarget{data-gaps-and-ongoing-work}{%
\subsubsection{5.8 Data Gaps and Ongoing
Work}\label{data-gaps-and-ongoing-work}}

Table 6 consolidates the coverage gaps across all experimental
configurations. Five of seven configurations are complete (3 seeds, 200
epochs, all 7 methods); the remaining two have partial data.

\textbf{Table 6: Data gap summary. Complete = 7 methods \(\times\) 3
seeds \(\times\) 200 epochs.}

\begin{longtable}[]{@{}
  >{\raggedright\arraybackslash}p{(\columnwidth - 6\tabcolsep) * \real{0.2373}}
  >{\raggedright\arraybackslash}p{(\columnwidth - 6\tabcolsep) * \real{0.1356}}
  >{\raggedright\arraybackslash}p{(\columnwidth - 6\tabcolsep) * \real{0.2542}}
  >{\raggedright\arraybackslash}p{(\columnwidth - 6\tabcolsep) * \real{0.3729}}@{}}
\toprule\noalign{}
\begin{minipage}[b]{\linewidth}\raggedright
Configuration
\end{minipage} & \begin{minipage}[b]{\linewidth}\raggedright
Status
\end{minipage} & \begin{minipage}[b]{\linewidth}\raggedright
Available Data
\end{minipage} & \begin{minipage}[b]{\linewidth}\raggedright
What It Would Resolve
\end{minipage} \\
\midrule\noalign{}
\endhead
\bottomrule\noalign{}
\endlastfoot
\(\rho\)-high (AdamW, \(\rho \approx 5.0\)) & Pilot only & 5-epoch
pilot: 77.69\% & Theorem 1 prediction at elevated \(\rho\) \\
Matched-\(\rho\) SGD & Partial & Constant: 2 seeds (200 ep); CWD: 1 seed
(200 ep) & Optimizer-vs-\(\rho\) confound resolution \\
NoBN & Partial & 2/7 methods complete (3 seeds each) & Full NoBN Phi
spread; Theorem 1 threshold test \\
\(\rho\)-low (AdamW, \(\rho \approx 0.05\)) & Partial & Constant: 3
seeds; CWD: 1 seed & Low-\(\rho\) regime Phi spread \\
\end{longtable}

All 5 complete configurations confirm Theorem 1's predictions. The
partial data available from incomplete configurations is directionally
consistent but insufficient for statistical conclusions.

\begin{center}\rule{0.5\linewidth}{0.5pt}\end{center}

\hypertarget{discussion}{%
\subsection{6. Discussion}\label{discussion}}

\hypertarget{why-constant-wd-wins-at-standard-rho}{%
\subsubsection{\texorpdfstring{6.1 Why Constant WD Wins at Standard
\(\rho\)}{6.1 Why Constant WD Wins at Standard \textbackslash rho}}\label{why-constant-wd-wins-at-standard-rho}}

Theorem 1 predicts this outcome through a quantitative tradeoff. At
standard \(\rho \approx 0.5\) under AdamW with
\(\lambda = 5 \times 10^{-4}\), BN networks satisfy two conditions that
make constant WD optimal:

\begin{enumerate}
\def\labelenumi{\arabic{enumi}.}
\item
  \textbf{BN's scale-invariance limits alignment benefit.} BN layers are
  invariant to weight scaling:
  \(\text{BN}(\alpha \theta) = \text{BN}(\theta)\). This drives weights
  toward an equilibrium norm (Kosson et al., 2023), making the
  gradient-weight alignment signal \(\hat{\delta}_t\) geometrically
  constrained rather than freely informative. AIS values in our BN
  experiments (0.18--0.40) fall below the Theorem 1 threshold.
\item
  \textbf{AdamW's per-parameter scaling subsumes \(\phi\) effects.}
  AdamW's second-moment normalization (\(\hat{v}_t^{-1/2}\)) already
  provides per-parameter adaptive scaling. The additional modulation
  from \(\phi\) is redundant: it adds noise to an already
  well-calibrated update without providing new information. Across all 7
  methods, a 10\(\times\) variation in effective WD budget (BEM from 0.0
  to 0.90) produces \(< 0.25\%\) accuracy change.
\end{enumerate}

\hypertarget{when-dynamic-wd-should-help}{%
\subsubsection{6.2 When Dynamic WD Should
Help}\label{when-dynamic-wd-should-help}}

The theory predicts three regimes where dynamic WD becomes beneficial,
each with different levels of empirical support:

\textbf{Elevated \(\rho\) (partially supported; Section 5.3).} When
\(\rho > \rho^*\), the alignment benefit exceeds the stability cost
(Theorem 1). The \(\rho\)-low partial data
(\(\Phi_{\text{spread}} \geq 0.04\%\)) and SGD data
(\(\Phi_{\text{spread}} = 0.91\%\)) provide two points on the
spread-vs-\(\rho\) curve, both directionally consistent. Full
\(\rho\)-high data (\(\rho \approx 5.0\)) is needed to confirm the
predicted increase at extreme ratios (Table 6).

\textbf{Without BN (directional; Section 5.5).} NoBN experiments show
higher AIS (\$\sim\$0.50 vs.~\$\sim\$0.35), consistent with the
prediction that removing BN's scale-invariance increases alignment
informativeness. The available NoBN data (2 methods, constant
outperforming CWD by 0.10\%, within 1-std margin) does not yet resolve
whether the AIS increase is sufficient to cross the Theorem 1 threshold;
full NoBN Phi spread requires additional methods (Table 6).

\textbf{Large-batch training (speculative; no experimental data).}
Proposition 1's noise constraint relaxes as batch size increases:
\(\sigma^2/n\) decreases, lowering the Theorem 1 threshold. At LLM-scale
batch sizes (4K--64K tokens), the raw alignment signal becomes more
informative, potentially enabling direct alignment-based WD modulation
without EMA smoothing. This remains unvalidated.

\hypertarget{pmp-wd-as-a-principled-alternative}{%
\subsubsection{6.3 PMP-WD as a Principled
Alternative}\label{pmp-wd-as-a-principled-alternative}}

PMP-WD is derived from an optimality condition on the training dynamics.
Three properties distinguish it from existing methods:

\begin{enumerate}
\def\labelenumi{\arabic{enumi}.}
\item
  \textbf{State-feedback vs.~feedforward (Section 3.5, Figure 3).}
  PMP-WD is designed to correct deviations from \(\rho^*\) in real time;
  feedforward methods (cosine, AdamC) cannot compensate when the actual
  trajectory deviates from plan. Empirical validation is pending
  (Section 5.8).
\item
  \textbf{Dual derivation.} The same functional form emerges from both
  stochastic PMP (control theory) and RG beta function analysis
  (statistical physics), with agreement within 15\% over the
  moderate-alignment regime \(\hat{\delta}_t \in [0.3, 0.7]\) (Remark
  3.1).
\item
  \textbf{Minimal overhead.} PMP-WD requires only per-layer scalar EMA
  of \(\|g_l\| / \|\theta_l\|\)---a computation already performed for
  gradient norm tracking in standard training loops. The additional cost
  is \(O(L)\) per step, negligible relative to the forward/backward
  pass.
\end{enumerate}

\hypertarget{implications-for-practitioners}{%
\subsubsection{6.4 Implications for
Practitioners}\label{implications-for-practitioners}}

Three practical recommendations follow from the theory and experiments:

\begin{enumerate}
\def\labelenumi{\arabic{enumi}.}
\item
  \textbf{Under AdamW at standard hyperparameters
  (\(\lambda \sim 10^{-4}\)--\(10^{-3}\), BN architectures): use
  constant WD.} Dynamic scheduling adds implementation complexity
  without accuracy benefit. The 0.25\% Phi spread on CIFAR-10 is within
  the 0.31\% inter-seed std of the best method; VGG-16-BN's 0.16\%
  spread falls below any single method's confidence interval width.
\item
  \textbf{New dynamic WD methods should demonstrate gains at elevated
  \(\rho\) or at scale.} Standard CIFAR settings with AdamW are in the
  Phi-invariance regime; positive results here do not generalize to the
  claim that the method is universally better. The gradient-to-weight
  ratio should be reported as a context variable.
\item
  \textbf{If using alignment feedback, EMA smoothing with \(k \geq 10\)
  is mandatory.} Proposition 1 shows that raw single-step alignment
  signals are unreliable at batch sizes \(\leq\) 256. CWD's binary mask
  is partially robust (sign is low-dimensional), but continuous
  alignment methods must aggregate temporally.
\end{enumerate}

\hypertarget{limitations}{%
\subsubsection{6.5 Limitations}\label{limitations}}

\begin{enumerate}
\def\labelenumi{\arabic{enumi}.}
\item
  \textbf{Scale.} CIFAR-10/100 with ResNet-20 (270K parameters) and
  VGG-16-BN (15M parameters) are small by current standards. ImageNet
  and LLM experiments are not included. The theory's predictions at
  larger scale remain unvalidated.
\item
  \textbf{Incomplete configurations.} Three configurations
  (\(\rho\)-high, matched-\(\rho\) SGD, NoBN) have partial data; Table 6
  (Section 5.8) details the specific gaps and what each would resolve.
\item
  \textbf{PMP-WD not implemented.} The algorithm is derived but not
  empirically evaluated. Its predicted advantage at elevated \(\rho\) is
  theoretical.
\item
  \textbf{Three seeds.} With 3 seeds, statistical power is limited for
  effect sizes below 0.3\%. TOST equivalence testing at margin
  \(\pm 0.3\%\) requires larger \(n\) for definitive equivalence claims.
\item
  \textbf{Fixed hyperparameters.} All methods share the same base
  \(\lambda\) and lr. Per-method hyperparameter tuning (e.g., CWD's
  \(\beta\) parameter) might reveal larger differences, at the cost of
  confounding method identity with hyperparameter optimization quality.
\end{enumerate}

\begin{center}\rule{0.5\linewidth}{0.5pt}\end{center}

\hypertarget{conclusion}{%
\subsection{7. Conclusion}\label{conclusion}}

Dynamic WD methods fail to outperform constant WD at standard settings
not by accident but by necessity: Theorem 1 shows that the alignment
benefit of \(\phi\) modulation is exceeded by its stability cost
whenever the gradient-to-weight ratio \(\rho\) falls below a
noise-dependent threshold, a condition satisfied in all 5 complete
configurations tested (105 completed 200-epoch runs). Theorem 2 explains
why binary masking methods (CWD, random mask) incur hidden per-parameter
stability costs despite moderate aggregate Coupling Stability Index
(CSI), by bounding the generalization penalty under per-parameter WD
variation. Together, Theorems 1--2 and Proposition 1 (requiring EMA
aggregation \(k \geq 10\) for reliable alignment feedback) provide a
quantitative theory of when and why constant WD is optimal. Theorem 3
derives a state-feedback WD controller (PMP-WD) from the stochastic
Pontryagin Maximum Principle and independently from renormalization
group beta function theory; this algorithm is derived but not yet
empirically evaluated, with validation at elevated \(\rho\) as a
priority for future work.

SGD shows 3.7\(\times\) larger method sensitivity than AdamW on CIFAR-10
(0.91\% vs.~0.25\%), consistent with Theorem 1's prediction that Phi
spread scales with \(\rho\), though direct matched-\(\rho\) comparison
remains to be completed (Section 6.5). The practical implication is
direct: under AdamW at standard hyperparameters, constant WD is
sufficient; dynamic WD methods should target elevated-\(\rho\) or
large-scale regimes where the Alignment Informativeness Score (AIS)
exceeds the stability cost.

Immediate extensions include ImageNet-scale validation and PMP-WD
empirical evaluation at elevated \(\rho\), which directly tests Theorem
3's predictions. Longer-horizon directions include Vision Transformer
architectures, where the absence of BN may increase AIS (Section 6.2),
and LLM-scale experiments, where large batch sizes relax Proposition 1's
noise floor and long training horizons may amplify WD timing effects.

\begin{center}\rule{0.5\linewidth}{0.5pt}\end{center}

\hypertarget{figures-and-tables}{%
\subsection{Figures and Tables}\label{figures-and-tables}}

\begin{itemize}
\tightlist
\item
  Figure 1: ratio\_regime.pdf --- Ratio regime diagram: method spread
  vs.~log(\(\rho\)) with three regime zones
\item
  Figure 2: theorem1\_regime.pdf --- Theorem 1 regime illustration:
  alignment benefit vs.~stability cost crossover at AIS*
\item
  Figure 3: pmpwd\_control.pdf --- PMP-WD control diagram: closed-loop
  state-feedback vs.~feedforward
\item
  Figure 4: multi\_arch\_comparison.pdf --- Multi-architecture accuracy
  comparison (ResNet-20 and VGG-16-BN)
\item
  Figure 5: sgd\_vs\_adamw\_spread.pdf --- Phi spread comparison across
  optimizers
\item
  Figure 6: diagnostic\_panel.pdf --- Diagnostic metric panel (CSI, AIS,
  BEM vs.~accuracy)
\item
  Table 1: inline --- Main accuracy results (7 methods \(\times\) 4
  configurations)
\item
  Table 2: inline --- Statistical significance tests (AdamW CIFAR-10)
\item
  Table 3: inline --- VGG-16-BN results
\item
  Table 4: inline --- Method taxonomy (Phi modulator special cases)
\item
  Table 5: inline --- NoBN vs.~BN ablation
\item
  Table 6: inline --- Data gap summary
\end{itemize}

\end{document}
